% GENERATED, edit the French unit or the translation table
\mxcorrections{5}\label{C05-CORR}
\textbf{Teacher use.} The answers below concern the tasks not worked in the lesson. First have students state their method; give graded help before the solution. The durations are estimates to be calibrated in class.
\subsection*{Diagnostic and remediation}
\begin{enumerate}
\item $\dfrac{12}{40}=\dfrac{3}{10}=30\,\%$. \textit{Action:} Link the partial count to the reference count.
\item $0.7$. \textit{Action:} Represent the shared sample space as two events.
\item $A\cap B$. \textit{Action:} Shade the intersection of two sets.
\item $A\cup B$. \textit{Action:} Include the common outcomes in the union.
\item $0.1$. \textit{Action:} Replace the decimals by fractions.
\item $\dfrac{3}{6}=\dfrac{1}{2}$. \textit{Action:} Identify the equally-likely assumption.
\end{enumerate}
\subsection*{Investigation activity}
Required intersection 12; cells 12, 28, 18, 42. Replacing them by 13, 27, 17, 43 keeps the totals but $\dfrac{13}{100}\ne 0.4\times 0.3$.
\subsection*{Immediate applications}
\begin{enumerate}
\item $\dfrac{5}{60}=\dfrac{1}{12}$ and $\dfrac{5}{20}=\dfrac{1}{4}$.
\item $\dfrac{1}{4}+\dfrac{1}{6}=\dfrac{5}{12}$.
\item No: $0.4\times0.5=0.20\ne0.18$.
\item $(\dfrac{2}{3})^4=\dfrac{16}{81}$ then $1-\dfrac{16}{81}=\dfrac{65}{81}$.
\item $\dfrac{675}{2000}=0.3375$; this is an estimate, not the exact value.
\end{enumerate}
\subsection*{Exercises: method and results}
\paragraph{Exercise 5.1 -- A table.} Cells: 15 (A and B), 25 (A without B), 20 (B without A), 40 (neither A nor B). $P_A(B)=\dfrac{15}{40}=\dfrac{3}{8}$; $P_B(A)=\dfrac{15}{35}=\dfrac{3}{7}$.
\paragraph{Exercise 5.2 -- Production.} $P(D)=0.25\times0.04+0.75\times0.02=0.025$. $P_D(A)=\dfrac{0.01}{0.025}=0.4$. The denominator changes when we condition on D.
\paragraph{Exercise 5.3 -- Checking independence.} $P(A)=P(B)=\dfrac{1}{2}$; the intersection $\{4;6\}$ has probability $\dfrac{1}{3}\ne \dfrac{1}{4}$: not independent. The fact that another choice of threshold gave independence does not transfer automatically.
\paragraph{Exercise 5.4 -- Without replacement.} Two red: $(\dfrac{3}{5})(\dfrac{2}{4})=\dfrac{3}{10}$. One of each: $(\dfrac{3}{5})(\dfrac{2}{4})+(\dfrac{2}{5})(\dfrac{3}{4})=\dfrac{3}{5}$. The order distinguishes two disjoint paths.
\paragraph{Exercise 5.5 -- Three trials.} One success: three paths, each $(\dfrac{1}{3})(\dfrac{2}{3})^2=\dfrac{4}{27}$, total $\dfrac{4}{9}$. Two successes: $3(\dfrac{1}{3})^2(\dfrac{2}{3})=\dfrac{2}{9}$; three: $\dfrac{1}{27}$; total $\dfrac{7}{27}$.
\paragraph{Exercise 5.6 -- Detector.} $P(S)=0.02(0.9)+0.98(0.05)=0.067$. $P_S(D)=\dfrac{0.018}{0.067}=\dfrac{18}{67}\approx 26.9\,\%$. The low initial rate explains the gap with 90 \%.
\paragraph{Exercise 5.7 -- Sets and cardinality.} $A\cap B=\{2\}$ (1); $A\cup B=\{1;2;4\}$ (3); $E\setminus A=\{3;4\}$ (2); product $\{(1;2),(1;4),(2;2),(2;4)\}$ (4). A pair is ordered.
\paragraph{Exercise 5.8 -- Delayed review.} $P(B)=\dfrac{1}{8}+\dfrac{3}{8}=\dfrac{1}{2}$. $P_B(A)=\dfrac{\dfrac{1}{8}}{\dfrac{1}{2}}=\dfrac{1}{4}$. The numerical equality with $P_A(B)$ is particular to these data.
\paragraph{CQFD, exercise 67 p.~336, questions 1, 2 and 3.} $P(\overline D\cap\overline N)=0.9\times0.05=0.045$. The probability of a positive test is $0.1\times0.97+0.045=0.142$. Among the positive tests, the probability of not being doped is $\dfrac{0.045}{0.142}=\dfrac{45}{142}\approx31.7\,\%$. Do not swap positive and negative test.
\paragraph{CQFD, exercise 70 p.~336, questions a and b.} $P(B)=\dfrac14$, $P_A(B)=\dfrac12$ and $P_{\overline A}(B)=0$. Hence $\dfrac14\ne\dfrac12+0$, then $0\ne1-\dfrac12$. The law of total probability involves the weights $P(A)$ and $P(\overline A)$; the complement applies to B when the same condition is held fixed.
\subsection*{MCQ: answers and reasons}
\textbf{MCQ 1.}
\begin{enumerate}
\item B. Reference sample space A.
\item A. Product of the branches.
\item D. Definition.
\item C. Criterion without division.
\item B. Zero intersection, positive product.
\end{enumerate}
\textbf{MCQ 2.}
\begin{enumerate}
\item C. Complement of three faces.
\item A. Product $p^2$.
\item D. Probabilities of a conditional partition.
\item B. Observed frequency and model probability differ.
\item C. Two red balls remain out of 4.
\end{enumerate}
\subsection*{True or false}
\begin{enumerate}
\item False: 12 common objects out of 40 A and 30 B give $\dfrac{3}{10}$ and $\dfrac{2}{5}$.
\item True: $P(A\cap\overline B)=P(A)-P(A\cap B)=P(A)(1-P(B))$.
\item False: without replacement, the composition of the urn changes.
\end{enumerate}
\paragraph{Differentiated teaching.} For a technical block, offer only the first step (domain, factor or reference sample space), then withdraw the help. For an independent student, ask for a second method and the treatment of a limiting case. Orally, allow an initial reformulation in the language the student has mastered, then require a precise mathematical sentence in French.
